% =====================================================================
% Artigo de TCC -- Bacharelado em Ciência da Computação -- IC/UFF
%
% Estrutura e formatação seguem o modelo oficial do curso
% (modelo_artigo/ModeloTCC_Artigo_CC_Latex.zip): capa, folha de aprovação,
% ficha catalográfica, título bilíngue, resumo/abstract, Introdução,
% Desenvolvimento, Conclusão, Referências, Apêndices e Agradecimentos.
%
% Overleaf: compilador pdfLaTeX, documento principal main.tex.
% Trechos marcados com \pendente{...} aparecem em vermelho no PDF e
% precisam ser completados ou removidos antes da versão final.
% =====================================================================
\documentclass[12pt,a4paper]{article}

% --- Codificação, idioma e fontes ------------------------------------
\usepackage[T1]{fontenc}
\usepackage[utf8]{inputenc}
\usepackage[brazil]{babel}
\usepackage{lmodern}
\usepackage{microtype}

% --- Página: A4, margens de 2,54 cm, espaçamento 1,5 e recuo de 1,25 cm
\usepackage[a4paper,margin=2.54cm]{geometry}
\usepackage{setspace}
\onehalfspacing
\usepackage{indentfirst}
\setlength{\parindent}{1.25cm}

% --- Matemática, tabelas e ilustrações --------------------------------
\usepackage{amsmath,amssymb}
\usepackage{array}
\usepackage{booktabs}
\usepackage{graphicx}
\usepackage{xcolor}
\usepackage{tikz}
\usetikzlibrary{arrows.meta,positioning}
\usepackage{pgfplots}
\pgfplotsset{compat=1.18}
\usepackage[ruled,vlined,linesnumbered,portuguese]{algorithm2e}
% Palavras-chave dos algoritmos em português.
\SetKwInput{KwIn}{Entrada}
\SetKwInput{KwOut}{Saída}
\SetKwFor{While}{enquanto}{faça}{fim}
\SetKwIF{If}{ElseIf}{Else}{se}{então}{senão se}{senão}{fim}
\SetKw{Return}{retorne}

% --- Legendas: "Figura 1 – ...", legenda de tabela acima -------------
\usepackage[labelsep=endash]{caption}
\captionsetup[table]{position=top}

% --- Títulos de seção como no modelo: "1. INTRODUÇÃO" ----------------
\usepackage{titlesec}
\titleformat{\section}{\normalfont\normalsize\bfseries}{\thesection.}{0.5em}{\MakeUppercase}
\titleformat{\subsection}{\normalfont\normalsize\bfseries}{\thesubsection.}{0.5em}{\MakeUppercase}
\titleformat{\subsubsection}{\normalfont\normalsize\bfseries}{\thesubsubsection.}{0.5em}{}
\titlespacing*{\section}{0pt}{1.5\baselineskip}{0.5\baselineskip}
\titlespacing*{\subsection}{0pt}{1.2\baselineskip}{0.5\baselineskip}
\titlespacing*{\subsubsection}{0pt}{\baselineskip}{0.3\baselineskip}

% --- Listas de ilustrações e tabelas (Apêndices A e B) ---------------
\usepackage[titles]{tocloft}
\renewcommand{\cftfigpresnum}{Figura~}
\renewcommand{\cftfigaftersnum}{~--}
\setlength{\cftfignumwidth}{2.7cm}
\setlength{\cftfigindent}{0pt}
\renewcommand{\cfttabpresnum}{Tabela~}
\renewcommand{\cfttabaftersnum}{~--}
\setlength{\cfttabnumwidth}{2.7cm}
\setlength{\cfttabindent}{0pt}

% --- Links e citações autor-data (NBR 10520) / referências (NBR 6023) -
% O hyperref precisa ser carregado ANTES do abntex2cite.
\usepackage{hyperref}
\usepackage[alf]{abntex2cite}
\usepackage{etoolbox}
\hypersetup{
  hidelinks,
  pdftitle={Alocação de professores, horários e salas no Instituto de Computação da UFF},
  pdfauthor={Bruno Costa dos Passos}
}

% --- Dados do trabalho -----------------------------------------------
\newcommand{\autor}{Bruno Costa dos Passos}
\newcommand{\orientador}{Igor Machado Coelho}
\newcommand{\titulo}{Alocação de professores, horários e salas no Instituto de Computação da UFF: uma abordagem meta-heurística com PyOptFrame}
\newcommand{\tituloingles}{Teacher, timeslot and room assignment at the Institute of Computing of UFF: a metaheuristic approach with PyOptFrame}
\newcommand{\anoentrega}{2026}
\newcommand{\natureza}{Trabalho de Conclusão de Curso submetido ao Departamento de Ciência da Computação da Universidade Federal Fluminense como requisito parcial para a obtenção do título de Bacharel em Ciência da Computação.}

% --- Utilitários ------------------------------------------------------
% Marcador de pendência: remover todas as ocorrências na versão final.
\newcommand{\pendente}[1]{\textcolor{red}{\textbf{[PENDENTE: #1]}}}
% Títulos pós-textuais (Referências, Apêndices, Agradecimentos).
\newcommand{\titulopostextual}[1]{%
  \clearpage
  \phantomsection
  \begin{center}\LARGE #1\end{center}
  \vspace{\baselineskip}}
% Referências com o mesmo título centralizado do modelo.
\makeatletter
\patchcmd{\thebibliographyBkUp}{\section*{\refname}}{\titulopostextual{\refname}}{}%
  {\PackageWarningNoLine{main}{Nao foi possivel ajustar o titulo das Referencias}}
\makeatother

\begin{document}

% =====================================================================
% ELEMENTOS PRÉ-TEXTUAIS
% =====================================================================
% =====================================================================
% Elementos pré-textuais, na ordem do modelo do curso:
% capa (com a natureza do trabalho), folha de aprovação e ficha
% catalográfica. Essas páginas contam na numeração, mas não a exibem;
% por isso o artigo começa na página 4, como no modelo.
% =====================================================================

% ---------------------------------------------------------------------
% Capa
% ---------------------------------------------------------------------
\thispagestyle{empty}
\begin{center}
  \begin{singlespace}
    UNIVERSIDADE FEDERAL FLUMINENSE\\
    INSTITUTO DE COMPUTAÇÃO\\
    DEPARTAMENTO DE CIÊNCIA DA COMPUTAÇÃO
  \end{singlespace}

  \vspace{2.6cm}
  \MakeUppercase{\autor}

  \vspace{5.2cm}
  \begin{singlespace}
    \MakeUppercase{\titulo}
  \end{singlespace}
\end{center}

\vspace{2.2cm}
\hfill
\begin{minipage}{0.55\textwidth}
  \singlespacing
  \natureza
\end{minipage}

\vfill
\begin{center}
  {\large Niterói\\ \anoentrega}
\end{center}
\clearpage

% ---------------------------------------------------------------------
% Folha de aprovação
% ---------------------------------------------------------------------
\thispagestyle{empty}
\begin{center}
  \MakeUppercase{\autor}

  \vspace{3.8cm}
  \begin{singlespace}
    \MakeUppercase{\titulo}
  \end{singlespace}
\end{center}

\vspace{1.4cm}
\hfill
\begin{minipage}{0.55\textwidth}
  \singlespacing
  \natureza
\end{minipage}

\vspace{1.6cm}
\hspace*{1.25cm}Aprovado em \pendente{XX de XXX de XXXX}:

\vspace{1.2cm}
\begin{singlespace}
\begin{center}
  \rule{9cm}{0.4pt}\\
  \textbf{Prof. Dr. \orientador}\\
  Orientador -- UFF

  \vspace{1.3cm}
  \rule{9cm}{0.4pt}\\
  \textbf{\pendente{nome do membro da banca}}\\
  \pendente{instituição}

  \vspace{1.3cm}
  \rule{9cm}{0.4pt}\\
  \textbf{\pendente{nome do membro da banca}}\\
  \pendente{instituição}
\end{center}
\end{singlespace}

\vfill
\begin{center}
  {\large Niterói\\ \anoentrega}
\end{center}
\clearpage

% ---------------------------------------------------------------------
% Ficha catalográfica
% Gere a ficha em https://bibliotecas.uff.br/bee/fichacatalografica,
% salve o PDF como "ficha_catalografica.pdf" nesta pasta e troque o bloco
% abaixo por:
%   \noindent\includegraphics[width=\textwidth]{ficha_catalografica.pdf}
% ---------------------------------------------------------------------
\thispagestyle{empty}
\vspace*{\fill}
\begin{center}
  \begin{singlespace}
    \small
    Inserir aqui a ficha gerada a partir do Sistema de Geração Automática de Fichas Catalográficas,\\
    disponível no endereço \url{https://bibliotecas.uff.br/bee/fichacatalografica}.\\[0.3cm]
    \pendente{gerar a ficha catalográfica para a versão final}
  \end{singlespace}
\end{center}
\clearpage

% A partir daqui as páginas exibem numeração (o artigo começa na página 4).
\pagestyle{plain}


% =====================================================================
% PRIMEIRA PÁGINA DO ARTIGO: títulos, autoria, resumo e abstract
% =====================================================================
% =====================================================================
% Primeira página do artigo: títulos, autoria, resumo e abstract.
% O resumo deve ser um parágrafo único com objetivo, método, resultados
% e conclusões. Atualizar após os experimentos definitivos.
% =====================================================================

\begin{center}
  \begin{singlespace}
    {\LARGE \titulo\par}
    \vspace{0.35cm}
    {\LARGE\itshape \tituloingles\par}
  \end{singlespace}
\end{center}

\vspace{0.4cm}
\begin{flushright}
  \autor\footnote{Graduando do Curso de Ciência da Computação -- UFF, e-mail: \pendente{e-mail id.uff do autor}.}\\
  \orientador\footnote{Orientador -- Instituto de Computação -- UFF, e-mail: \pendente{e-mail do orientador}.}
\end{flushright}

\vspace{0.4cm}
\begin{center}
  \textbf{Resumo}
\end{center}

\noindent A elaboração anual de quadros de horários universitários exige decidir, de forma coordenada, quem ministra cada turma, em que horário ela ocorre e em que sala acontece cada encontro semanal. Este trabalho aborda esse problema no Instituto de Computação da Universidade Federal Fluminense (IC/UFF). O problema é modelado como uma variante do \textit{Curriculum-Based Course Timetabling}, com atribuição de professores, horizonte anual, grades de Ciência da Computação e Sistemas de Informação, turmas compartilhadas e regras locais, como a carga mínima anual de três disciplinas obrigatórias por docente. Uma pipeline reprodutível extrai, audita e integra os quadros de horários de 2026 e as grades curriculares, resultando em uma instância com 180 turmas, 329 encontros semanais e 22 salas. Sobre essa base, foram implementados um avaliador decomposto por critério, movimentos de professor, horário e sala, uma melhoria gulosa e as meta-heurísticas \textit{Simulated Annealing} (SA), \textit{Iterated Local Search} e \textit{Variable Neighborhood Search}, além da integração com o SA nativo do pyoptframe. Em um perfil sintético que completa os dados institucionais ainda não validados, a melhor execução do SA reduziu de 27 para 10 as violações de restrições fortes do quadro observado. Em um cenário exploratório com horários flexíveis para as obrigatórias, a redução chegou a zero, ao custo de mais janelas e dias de aula para os docentes. Os resultados indicam que a abordagem é viável, mas ainda dependem da validação dos dados e das hipóteses institucionais. \pendente{atualizar com os resultados da instância definitiva.}

\vspace{0.5\baselineskip}
\noindent\textbf{Palavras-chave:} Timetabling universitário; Meta-heurísticas; Atribuição de professores; Alocação de salas; OptFrame.

\vspace{\baselineskip}
\begin{center}
  \textbf{Abstract}
\end{center}

\noindent Building a university timetable every year requires deciding, in a coordinated way, who teaches each class, when it takes place, and in which room each weekly meeting happens. This work addresses this problem at the Institute of Computing of Fluminense Federal University (IC/UFF). The problem is modeled as a variant of Curriculum-Based Course Timetabling that includes teacher assignment, an annual planning horizon, the Computer Science and Information Systems curricula, shared classes, and local rules such as a minimum annual load of three mandatory courses per teacher. A reproducible pipeline extracts, audits, and integrates the 2026 timetables and curricula, yielding an instance with 180 classes, 329 weekly meetings, and 22 rooms. On top of this data, the work implements an evaluator decomposed by criterion, teacher, timeslot, and room moves, a greedy improvement step, and the Simulated Annealing (SA), Iterated Local Search, and Variable Neighborhood Search metaheuristics, together with an integration with the native SA of pyoptframe. On a synthetic profile that completes the institutional data still under validation, the best SA run reduced the hard constraint violations of the observed timetable from 27 to 10. In an exploratory scenario with flexible schedules for mandatory courses, the reduction reached zero, at the cost of more idle periods and teaching days for the teachers. The results indicate that the approach is feasible, but they still depend on the validation of institutional data and assumptions. \pendente{update with the results of the final instance.}

\vspace{0.5\baselineskip}
\noindent\textbf{Keywords:} University course timetabling; Metaheuristics; Teacher assignment; Room assignment; OptFrame.

\vspace{\baselineskip}
\noindent\textbf{Aprovado em: dd/mm/aaaa. Versão Final em: dd/mm/aaaa}


% =====================================================================
% ELEMENTOS TEXTUAIS
% =====================================================================
% =====================================================================
% 1. INTRODUÇÃO
% =====================================================================
\section{Introdução}
\label{sec:introducao}

A elaboração de quadros de horários e a alocação de salas em instituições de ensino superior são problemas clássicos de otimização combinatória, estudados na literatura de \textit{educational timetabling} \cite{schaerf1999survey}. No caso universitário, esse conjunto de problemas é usualmente chamado de \textit{University Course Timetabling Problem} (UCTP) \cite{lewis2008survey,babaei2015survey}. Mesmo versões simplificadas contêm problemas NP-difíceis como casos particulares. Além disso, essas versões combinam restrições de viabilidade com critérios de qualidade que frequentemente competem entre si.

No Instituto de Computação da Universidade Federal Fluminense (IC/UFF), a cada ano é preciso decidir, para cada turma ofertada nos dois semestres, qual professor a ministra, em que horário ela ocorre e em que sala acontece cada encontro semanal. Conforme as diretrizes levantadas com o orientador, essa decisão envolve três papéis institucionais. A chefia de departamento distribui os encargos docentes. A coordenação de curso representa os alunos, que precisam cursar as disciplinas previstas sem choques de horário e com vagas suficientes. O Instituto administra as salas e recebe pedidos de capacidade e de laboratórios. Esses papéis compartilham recursos, mas avaliam uma mesma grade sob perspectivas diferentes; turmas maiores, por exemplo, reduzem a demanda por professores, mas exigem salas maiores.

Algumas práticas locais tornam o problema mais estruturado. Muitas disciplinas seguem padrões de horário que se repetem a cada semestre ímpar ou par. As disciplinas se agrupam em setores do departamento, e cada setor tende a ocupar os mesmos dias da semana. Disciplinas oferecidas a outros cursos, como Estruturas de Dados, têm horário prioritário. Há ainda uma carga mínima anual de três disciplinas obrigatórias por professor, que acopla os dois semestres, e duas grades curriculares, Ciência da Computação (CC) e Sistemas de Informação (SI), que compartilham professores, salas e parte das turmas. \pendente{inserir a fonte normativa das regras de carga anual, prioridade docente e descanso entre jornadas, se houver documento institucional.}

A combinação dessas regras torna a construção manual da grade trabalhosa e difícil de auditar. Um modelo formal, acompanhado de métodos computacionais de busca, permite gerar alternativas e avaliar a grade praticada com os mesmos critérios aplicados às soluções propostas. O propósito deste trabalho, alinhado à orientação recebida, não é entregar um sistema pronto para uso administrativo. O foco é construir um modelo próximo da realidade observada, com hipóteses e origem dos dados documentadas, que sirva de base para trabalhos futuros. O critério de distância entre salas de aulas consecutivas, previsto na proposta inicial, foi retirado do escopo em comum acordo com o orientador.

\subsection{Objetivos}
\label{subsec:objetivos}

O objetivo geral deste trabalho é desenvolver e avaliar uma abordagem meta-heurística, implementada com o pyoptframe, para apoiar a alocação anual de professores, horários e salas no IC/UFF. A abordagem considera conflitos de horário, habilitação docente, jornada de trabalho, capacidade e recursos das salas e as preferências dos professores.

Os objetivos específicos são:
\begin{enumerate}
  \item formalizar o problema como uma variante do \textit{Curriculum-Based Course Timetabling} (CB-CTT), incorporando a atribuição de professores, o horizonte anual e as regras do IC/UFF;
  \item construir uma pipeline reprodutível de dados a partir dos quadros de horários de 2026, das páginas públicas de oferta e das grades de CC e SI, com auditoria e controle de prontidão da instância;
  \item implementar a representação da solução, um avaliador decomposto por critério, movimentos de professor, horário e sala e buscas integradas baseadas em meta-heurísticas, com integração ao pyoptframe;
  \item comparar as soluções geradas com a alocação observada, avaliada pelos mesmos critérios;
  \item estudar o efeito da ordem de planejamento das grades de CC e SI sobre a qualidade da solução. \pendente{confirmar se os cenários E1--E3 permanecem no escopo final.}
\end{enumerate}

\subsection{Contribuições}
\label{subsec:contribuicoes}

As principais contribuições obtidas até o momento são: (i) uma formulação do problema que integra professores, horários e salas por encontro em horizonte anual, com duas grades curriculares; (ii) uma pipeline que transforma os quadros de horários e as grades em uma instância auditável, separando dados observados, decisões humanas pendentes e hipóteses sintéticas; (iii) um solver com melhoria gulosa, \textit{Simulated Annealing} (SA), \textit{Iterated Local Search} (ILS) e \textit{Variable Neighborhood Search} (VNS), verificado por dois avaliadores e por uma validação estrutural; e (iv) uma primeira evidência empírica, em perfil sintético, de que a busca reduz substancialmente as violações de restrições fortes da grade observada. Também se identificou uma questão de modelagem relevante: a forma de tratar turmas alternativas na restrição de choque curricular.

\subsection{Organização do texto}
\label{subsec:organizacao}

A Seção~\ref{sec:desenvolvimento} apresenta o desenvolvimento do trabalho. Ela reúne a fundamentação teórica, a caracterização e a modelagem matemática do problema, os dados e as instâncias, o método de solução, o protocolo experimental e os resultados. A Seção~\ref{sec:conclusao} apresenta as conclusões, as limitações e os trabalhos futuros.


\section{Desenvolvimento}
\label{sec:desenvolvimento}

Esta seção reúne a fundamentação teórica (Seção~\ref{subsec:fundamentacao}), a caracterização e a modelagem do problema (Seções~\ref{subsec:problema} e \ref{subsec:modelagem}), os dados e as instâncias (Seção~\ref{subsec:dados}), o método de solução (Seção~\ref{subsec:metodo}), o protocolo experimental (Seção~\ref{subsec:protocolo}) e os resultados obtidos até o momento (Seção~\ref{subsec:resultados}).

% =====================================================================
% 2.1 FUNDAMENTAÇÃO TEÓRICA
% =====================================================================
\subsection{Fundamentação teórica}
\label{subsec:fundamentacao}

\subsubsection{Timetabling educacional}

O termo \textit{timetabling} designa a atribuição de eventos a períodos e, conforme a formulação, a outros recursos, sujeita a restrições de viabilidade e a critérios de qualidade. \citeonline{schaerf1999survey} organiza os problemas educacionais em três categorias: \textit{school timetabling}, que atribui aulas periódicas a professores e turmas escolares; \textit{examination timetabling}, que agenda provas cujos conflitos decorrem dos estudantes inscritos; e \textit{university course timetabling}, que agenda aulas universitárias considerando professores, grupos de estudantes, períodos e salas.

As restrições exatas variam de uma instituição para outra. Por isso, o UCTP designa uma família de problemas, e não uma formulação única \cite{lewis2008survey,babaei2015survey}. Em geral, distinguem-se restrições fortes (\textit{hard}), que definem a viabilidade de uma grade, e restrições fracas (\textit{soft}), que medem sua qualidade e costumam compor uma função objetivo ponderada.

\subsubsection{Curriculum-Based Course Timetabling}

Duas formulações universitárias recorrentes diferem pela origem dos conflitos entre estudantes. No \textit{Post-Enrolment Course Timetabling}, os conflitos derivam das matrículas individuais. No \textit{Curriculum-Based Course Timetabling} (CB-CTT), eles são definidos por \textit{curricula}, isto é, grupos de disciplinas que um mesmo conjunto de alunos deve cursar em paralelo.

A trilha 3 da \textit{Second International Timetabling Competition} (ITC-2007) consolidou uma formulação de referência para o CB-CTT \cite{digaspero2007itc}. Nela, uma solução atribui a cada aula um período e uma sala. Indisponibilidades, conflitos de professor ou de currículo e a ocupação exclusiva de sala são restrições fortes. Capacidade da sala, número mínimo de dias de trabalho, compactação curricular e estabilidade de sala aparecem como critérios fracos. \citeonline{bettinelli2015overview} apresentam uma visão geral das formulações e dos métodos propostos para o CB-CTT.

Este trabalho usa o CB-CTT como referência, mas não o reproduz literalmente. No IC/UFF, a atribuição do professor é uma decisão do modelo, enquanto no benchmark da ITC-2007 o professor de cada disciplina é dado. Além disso, o horizonte é anual, a sala é decidida por encontro semanal e há regras locais de setores, horários fixos e carga docente.

\subsubsection{Complexidade computacional}

Uma versão simplificada do problema de horários contém o problema de coloração de vértices. Considere um grafo de conflitos $G=(V,E)$ no qual cada vértice representa um evento e cada aresta indica que dois eventos não podem ocorrer simultaneamente. Associar um de $K$ períodos a cada evento, sem conflitos, equivale a obter uma $K$-coloração de $G$, isto é,
\begin{equation}
  c(u)\neq c(v), \qquad \forall (u,v)\in E.
\end{equation}
Como decidir a existência de uma $K$-coloração é NP-completo para $K\geq 3$ \cite{garey1979computers}, as formulações de otimização que generalizam essa versão são NP-difíceis. Restrições adicionais de salas, recursos e preferências não eliminam essa dificuldade, o que motiva o uso de métodos heurísticos em instâncias de porte realista.

\subsubsection{Meta-heurísticas de trajetória}

Meta-heurísticas exploram o espaço de soluções sem garantir otimalidade, buscando soluções de boa qualidade dentro de um orçamento computacional \cite{talbi2009metaheuristics}. Seja $\mathcal{S}$ o espaço de soluções e $F:\mathcal{S}\rightarrow\mathbb{R}$ uma função a minimizar. Uma estrutura de vizinhança $N$ associa a cada solução $s$ o conjunto $N(s)$ de soluções alcançáveis por um movimento elementar.

O \textit{Simulated Annealing} (SA), proposto por \citeonline{kirkpatrick1983optimization}, aceita ocasionalmente movimentos de piora para reduzir o risco de estagnação em ótimos locais. Dada uma solução corrente $s$ e uma vizinha $s'$, seja $\Delta=F(s')-F(s)$. Movimentos com $\Delta\leq 0$ são sempre aceitos; quando $\Delta>0$, a aceitação ocorre com probabilidade
\begin{equation}
  P(\text{aceitação})=\exp\left(-\frac{\Delta}{T}\right),
  \label{eq:metropolis}
\end{equation}
em que $T>0$ é a temperatura. A temperatura é reduzida ao longo da execução, o que torna a aceitação de pioras cada vez menos provável.

O \textit{Iterated Local Search} (ILS) alterna busca local, perturbação da melhor solução e um critério de aceitação, explorando uma sequência de ótimos locais \cite{lourenco2003iterated}. O \textit{Variable Neighborhood Search} (VNS) explora sistematicamente mais de uma estrutura de vizinhança, alternando entre elas para combinar intensificação e diversificação \cite{hansen2001variable}.

\subsubsection{OptFrame e pyoptframe}

O OptFrame é um framework para o desenvolvimento de métodos de otimização combinatória baseado em componentes reutilizáveis \cite{coelho2010optframe,coelho2020microbenchmark}. Entre esses componentes estão a representação da solução, o método construtivo, o avaliador, os movimentos, as estruturas de vizinhança e as próprias meta-heurísticas. O pyoptframe fornece \textit{bindings} Python para o núcleo funcional do OptFrame \cite{pyoptframe2026}.

A separação entre componentes permite trocar a meta-heurística ou a vizinhança sem reescrever o modelo do problema. Essa propriedade motivou a escolha do framework: representação, avaliação e movimentos puderam ser implementados e testados em Python e, depois, registrados no motor nativo.

\subsubsection{Trabalhos relacionados}
\label{subsubsec:relacionados}

Os levantamentos de \citeonline{lewis2008survey} e \citeonline{babaei2015survey} mostram que meta-heurísticas de trajetória e populacionais são amplamente empregadas em problemas de timetabling universitário. A visão geral de \citeonline{bettinelli2015overview} reúne formulações e abordagens para o CB-CTT, incluindo métodos exatos e heurísticos avaliados sobre as instâncias da ITC-2007.

\pendente{completar com trabalhos aplicados a instituições brasileiras (por exemplo, publicações do SBPO) e com estudos que integram atribuição de professores e alocação de salas. Para cada trabalho, indicar o problema tratado, o método e a diferença em relação a este TCC: horizonte anual, duas grades com turmas compartilhadas, sala por encontro e carga mínima anual de obrigatórias.}

% =====================================================================
% 2.2 CARACTERIZAÇÃO DO PROBLEMA e 2.3 MODELAGEM MATEMÁTICA
% =====================================================================
\subsection{Caracterização do problema no IC/UFF}
\label{subsec:problema}

A unidade programada é a \textit{turma}, isto é, a oferta de uma disciplina em um semestre. Cada turma possui um ou mais encontros semanais, cada um com dia, horário de início, horário de término e sala. A sala é decidida por encontro, e não por turma, porque há turmas que alternam uma aula teórica em sala comum e uma aula prática em laboratório. Para cada turma sob responsabilidade do IC, o modelo decide três aspectos acoplados: o professor, o padrão semanal de horário e a sala de cada encontro.

As turmas são classificadas por dois eixos independentes. Quanto à responsabilidade, distinguem-se as turmas do próprio IC, cujas decisões pertencem ao modelo, e as turmas de outros departamentos cursadas pelos alunos de CC e SI, como as de Cálculo. Estas últimas têm professor e horário dados e entram no modelo apenas como ocupação fixa do horário dos alunos. Quanto ao papel curricular, distinguem-se obrigatórias e optativas. Transversalmente, algumas turmas têm horário fixo: todas as de outros departamentos e as disciplinas de serviço do IC, oferecidas a outros cursos com prioridade.

O planejamento é anual. A carga mínima de obrigatórias e o rodízio de professores relacionam as decisões dos semestres ímpar e par, de modo que otimizá-los separadamente perderia essas relações. As restrições de choque, por outro lado, valem separadamente em cada semestre. O modelo considera ainda as grades de CC e SI, que pertencem ao mesmo departamento e seguem as mesmas regras. Uma disciplina presente nas duas grades corresponde a uma única turma, cursada por alunos dos dois cursos.

A Tabela~\ref{tab:papeis} relaciona os papéis institucionais e as regras locais às restrições e aos critérios do modelo apresentado na próxima seção.

\begin{table}[htbp]
  \centering
  \caption{Relação entre papéis institucionais, regras locais e componentes do modelo. Fonte: elaborada pelo autor (2026).}
  \label{tab:papeis}
  \small
  \begin{tabular}{p{6.2cm}p{8.2cm}}
    \toprule
    \textbf{Papel ou regra} & \textbf{Componentes do modelo} \\
    \midrule
    Chefia de departamento & Professor único (H1), carga anual (H12), preferência ponderada pela prioridade (O1) \\
    Coordenação de curso & Ausência de choque curricular (H8) e capacidade (H10) \\
    Instituto de Computação (salas) & Sala por encontro (H3), choque de sala (H7), recursos (H9), capacidade (H10) e desperdício (O4) \\
    Jornada do professor & Descanso entre jornadas (H11), dias com aula (O2) e janelas (O3) \\
    Disciplinas de serviço e de outros departamentos & Horário fixo (H4) e bloqueio do horário dos alunos (H8) \\
    Setores do departamento & Dias admissíveis por setor (H5) \\
    Grades de CC e SI & Grupos curriculares das duas grades em H8 \\
    \bottomrule
  \end{tabular}
\end{table}

\subsection{Modelagem matemática}
\label{subsec:modelagem}

\subsubsection{Conjuntos, parâmetros e variáveis}

Sejam $\mathcal{C}$ o conjunto de turmas do ano, particionado em turmas do IC ($\mathcal{C}^{\mathrm{IC}}$) e de outros departamentos ($\mathcal{C}^{\mathrm{out}}$), e também em obrigatórias ($\mathcal{C}^{\mathrm{obr}}$) e optativas ($\mathcal{C}^{\mathrm{opt}}$). O subconjunto $\mathcal{C}^{\mathrm{fix}}\supseteq\mathcal{C}^{\mathrm{out}}$ reúne as turmas com horário dado. O semestre de cada turma é $\sigma(c)\in\{\text{ímpar},\text{par}\}$, e $\mathcal{C}_\sigma$ denota as turmas do semestre $\sigma$. Os demais conjuntos são:
\begin{itemize}
  \setlength{\itemsep}{0pt}
  \item $\mathcal{P}$: professores do IC, com $\mathcal{P}^{\mathrm{H12}}\subseteq\mathcal{P}$ sujeitos à carga anual;
  \item $\mathcal{R}$: salas; $\mathcal{H}$: \textit{slots} semanais (dia $\times$ faixa horária);
  \item $\mathcal{Q}$: padrões semanais de horário, cada um com \textit{slots} $H_q\subseteq\mathcal{H}$;
  \item $\mathcal{M}_c$: encontros semanais da turma $c$;
  \item $\mathcal{G}$: grupos curriculares, cada um associado a uma grade, a um período e a um semestre, com turmas $\mathcal{C}_g\subseteq\mathcal{C}$.
\end{itemize}

Os domínios são $\mathcal{P}_c\subseteq\mathcal{P}$, professores habilitados para $c$; $\mathcal{Q}_c\subseteq\mathcal{Q}$, padrões admissíveis para $c$, unitário se $c\in\mathcal{C}^{\mathrm{fix}}$; e $\mathcal{R}_{c,m}\subseteq\mathcal{R}$, salas admissíveis para o encontro $m$. Os parâmetros principais são o número de vagas $n_c$, a capacidade $\operatorname{cap}_r$, os indicadores de exigência de laboratório $\rho_{c,m}$ e de sala laboratório $\ell_r$, a preferência $\operatorname{pref}_{p,c}\in[0,1]$, a prioridade $\pi_p\in[0,1]$, o padrão fixo $\bar q_c$ e o descanso mínimo entre jornadas $\operatorname{desc}_{\min}$. O parâmetro $a_{c,m,h,q}\in\{0,1\}$ indica se o encontro $m$ de $c$ ocupa o \textit{slot} $h$ quando o padrão $q$ é escolhido.

As variáveis de decisão são binárias:
\begin{align}
  x_{c,p}&=1 \text{ se a turma } c \text{ é atribuída ao professor } p, && c\in\mathcal{C}^{\mathrm{IC}},\ p\in\mathcal{P}_c,\\
  y_{c,q}&=1 \text{ se a turma } c \text{ recebe o padrão } q, && c\in\mathcal{C},\ q\in\mathcal{Q}_c,\\
  z_{c,m,r}&=1 \text{ se o encontro } m \text{ de } c \text{ ocorre na sala } r, && c\in\mathcal{C}^{\mathrm{IC}},\ m\in\mathcal{M}_c,\ r\in\mathcal{R}_{c,m}.
\end{align}
A ocupação de um \textit{slot} pela turma e pelo encontro é dada por
\begin{equation}
  u_{c,h}=\sum_{q\in\mathcal{Q}_c:\,h\in H_q} y_{c,q},
  \qquad
  v_{c,m,h}=\sum_{q\in\mathcal{Q}_c} a_{c,m,h,q}\,y_{c,q}.
\end{equation}

\subsubsection{Restrições fortes}

As restrições estruturais garantem um professor por turma do IC, um padrão de horário por turma, uma sala por encontro e o respeito aos horários fixos:
\begin{align}
  \sum_{p\in\mathcal{P}_c} x_{c,p} &= 1, && \forall c\in\mathcal{C}^{\mathrm{IC}}, \tag{H1}\\
  \sum_{q\in\mathcal{Q}_c} y_{c,q} &= 1, && \forall c\in\mathcal{C}, \tag{H2}\\
  \sum_{r\in\mathcal{R}_{c,m}} z_{c,m,r} &= 1, && \forall c\in\mathcal{C}^{\mathrm{IC}},\ m\in\mathcal{M}_c, \tag{H3}\\
  y_{c,\bar q_c} &= 1, && \forall c\in\mathcal{C}^{\mathrm{fix}}. \tag{H4}
\end{align}
A restrição de setor (H5) é imposta no domínio: para uma turma do IC com horário variável, $\mathcal{Q}_c$ contém apenas padrões cujos dias pertencem aos dias autorizados para o setor da disciplina.

As restrições de conflito valem para cada semestre $\sigma$. Um professor não ministra duas turmas no mesmo \textit{slot} (H6), uma sala não recebe dois encontros simultâneos (H7) e turmas de um mesmo grupo curricular não se sobrepõem (H8):
\begin{align}
  \sum_{c\in\mathcal{C}^{\mathrm{IC}}\cap\mathcal{C}_\sigma} x_{c,p}\,u_{c,h} &\leq 1, && \forall p\in\mathcal{P},\ h\in\mathcal{H},\ \sigma, \tag{H6}\\
  \sum_{c\in\mathcal{C}^{\mathrm{IC}}\cap\mathcal{C}_\sigma}\ \sum_{m\in\mathcal{M}_c} z_{c,m,r}\,v_{c,m,h} &\leq 1, && \forall r\in\mathcal{R},\ h\in\mathcal{H},\ \sigma, \tag{H7}\\
  \sum_{c\in\mathcal{C}_g} u_{c,h} &\leq 1, && \forall g\in\mathcal{G},\ h\in\mathcal{H}. \tag{H8}
\end{align}
Em H8, cada grupo contém apenas turmas do mesmo semestre, e as turmas de outros departamentos também participam, pois bloqueiam o horário dos alunos.

As restrições de recursos e capacidade (H9 e H10) restringem as salas admissíveis. A restrição de descanso (H11) proíbe que um professor lecione em dois \textit{slots} $(h,h')$ de dias consecutivos cujo intervalo seja menor que $\operatorname{desc}_{\min}$. Por fim, a carga anual (H12) exige ao menos três turmas obrigatórias do IC por professor do universo $\mathcal{P}^{\mathrm{H12}}$, somando os dois semestres:
\begin{align}
  z_{c,m,r} &= 0 \quad \text{se } \rho_{c,m}=1 \text{ e } \ell_r=0, \tag{H9}\\
  z_{c,m,r} &= 0 \quad \text{se } n_c>\operatorname{cap}_r, \tag{H10}\\
  \sum_{c\in\mathcal{C}^{\mathrm{IC}}\cap\mathcal{C}_\sigma} x_{c,p}\,u_{c,h}
  +\sum_{c\in\mathcal{C}^{\mathrm{IC}}\cap\mathcal{C}_\sigma} x_{c,p}\,u_{c,h'} &\leq 1,
  \quad \forall p,\ \sigma,\ (h,h')\in\mathcal{V}, \tag{H11}\\
  \sum_{c\in\mathcal{C}^{\mathrm{IC}}\cap\mathcal{C}^{\mathrm{obr}}} x_{c,p} &\geq 3,
  \quad \forall p\in\mathcal{P}^{\mathrm{H12}}, \tag{H12}
\end{align}
em que $\mathcal{V}$ é o conjunto de pares de \textit{slots} que violam o descanso mínimo. No protótipo, a compatibilidade de recursos é estrita: encontros sem exigência de laboratório também não ocupam laboratórios. As expressões H6, H7 e H11 contêm produtos de variáveis binárias apenas para apresentar a lógica de forma compacta; uma versão em programação linear inteira exigiria a linearização usual desses produtos.

\pendente{confirmar com o orientador a classificação final de H8, H10 e H11 como restrições fortes ou fracas, a composição de $\mathcal{P}^{\mathrm{H12}}$ e a fonte do descanso mínimo; a implementação atual usa $\operatorname{desc}_{\min}=11$ horas.}

\subsubsection{Função objetivo}

Os critérios de qualidade compõem a função objetivo
\begin{equation}
  \min Z = -w_{\mathrm{pref}}\Phi_{\mathrm{pref}}
  + w_{\mathrm{dias}}\Phi_{\mathrm{dias}}
  + w_{\mathrm{jan}}\Phi_{\mathrm{jan}}
  + w_{\mathrm{cap}}\Phi_{\mathrm{cap}}
  + w_{\mathrm{rod}}\Phi_{\mathrm{rod}},
  \label{eq:objetivo}
\end{equation}
cujos termos são:
\begin{itemize}
  \setlength{\itemsep}{0pt}
  \item (O1) $\Phi_{\mathrm{pref}}=\sum_{c,p}\pi_p\operatorname{pref}_{p,c}\,x_{c,p}$: preferências atendidas, ponderadas pela prioridade;
  \item (O2) $\Phi_{\mathrm{dias}}$: número de combinações professor--semestre--dia com pelo menos uma aula;
  \item (O3) $\Phi_{\mathrm{jan}}$: número de intervalos ociosos de duas horas entre aulas do mesmo professor no mesmo dia;
  \item (O4) $\Phi_{\mathrm{cap}}=\sum_{c,m,r} z_{c,m,r}(\operatorname{cap}_r-n_c)$: desperdício de capacidade;
  \item (O5) $\Phi_{\mathrm{rod}}$: número de pares professor--disciplina em que o professor ministra a mesma disciplina nos dois semestres do ano.
\end{itemize}

\subsubsection{Avaliação implementada}
\label{subsubsec:avaliacao}

A implementação não usa pesos numéricos para as restrições fortes. Ela compara soluções por uma chave lexicográfica
\begin{equation}
  \kappa(s)=\bigl(V_{\mathrm{hard}}(s),\ D_{\mathrm{H12}}(s),\ Z_1(s)\bigr),
  \label{eq:chave}
\end{equation}
em que $V_{\mathrm{hard}}$ soma sete contadores de violação (sala, professor, currículo, capacidade, recursos, descanso e docentes abaixo do mínimo anual) e $Z_1$ é a Equação~\eqref{eq:objetivo} com pesos unitários. O termo intermediário é o déficit de carga anual,
\begin{equation}
  D_{\mathrm{H12}}(s)=\sum_{p\in\mathcal{P}^{\mathrm{H12}}}\max\{0,\ 3-L_p(s)\},
\end{equation}
em que $L_p(s)$ é o número de obrigatórias do IC atribuídas a $p$ no ano. Turmas com dois professores contribuem com crédito fracionado para cada um. O déficit orienta a busca melhor que a contagem binária de docentes abaixo do mínimo, porque também diferencia soluções que ainda violam H12.

Os relatórios apresentam o \textit{score} $F(s)=10^6\,V_{\mathrm{hard}}(s)+Z_1(s)$. Para a regra de aceitação do SA, que exige um escalar, usa-se a energia $E(s)=10^9\,V_{\mathrm{hard}}(s)+10^6\,D_{\mathrm{H12}}(s)+Z_1(s)$, compatível com a ordem de $\kappa$. As restrições estruturais H1 a H5 são garantidas pela representação e pelos domínios dos movimentos, e uma validação estrutural independente as confere em toda solução retornada.

O contador de H8 merece uma observação. Ele conta sobreposições entre turmas de códigos distintos de um mesmo grupo e semestre, e ignora sobreposições entre seções de um mesmo código, que são alternativas entre si. Assim, duas seções alternativas de disciplinas diferentes que se sobrepõem contam como conflito, mesmo que o aluno possa escolher outra combinação de seções. As consequências dessa escolha são discutidas na Seção~\ref{subsubsec:h8}.

\pendente{definir com o orientador os pesos $w_\bullet$ dos critérios fracos; os experimentos atuais usam pesos unitários provisórios.}

% =====================================================================
% 2.4 DADOS E INSTÂNCIAS
% =====================================================================
\subsection{Dados e instâncias}
\label{subsec:dados}

\subsubsection{Fontes e pipeline de dados}

A instância foi construída a partir de quatro fontes. A primeira são os quadros de horários do IC de 2026/1 e 2026/2, em PDF, com turmas, professores, horários e salas. A segunda são as páginas públicas do Quadro de Horários da UFF, coletadas automaticamente, que informam vagas e inscritos por curso. A terceira são as grades curriculares de CC e SI, organizadas por período. A quarta é o histórico coletado de 2023/1 a 2025/2, junto com a planilha do quadro de 2025, usados como evidência indireta de preferências, setores e habilitações.

A Figura~\ref{fig:pipeline} resume o fluxo. A pipeline extrai e normaliza os registros, unifica grafias de nomes de docentes, vincula as linhas dos PDFs às páginas públicas e classifica cada oferta segundo as grades. Em seguida, audita a alocação observada e gera a instância em JSON. As decisões que dependem de validação humana, como capacidades físicas, setores e habilitações, são registradas em tabelas de revisão próprias, acompanhadas dos \textit{hashes} das fontes. Enquanto houver decisões pendentes, a instância institucional permanece marcada como não pronta para experimentos, e a pipeline não sobrescreve decisões humanas já registradas.

\begin{figure}[htbp]
  \centering
  \begin{tikzpicture}[
      node distance=0.45cm,
      caixa/.style={draw, rounded corners=2pt, align=center, text width=6.0cm,
                    minimum height=1.0cm, font=\small},
      seta/.style={-{Stealth[length=2.2mm]}, thick}]
    \node[caixa] (f) {Fontes: quadros de horários de 2026 (PDF), páginas públicas de oferta, grades de CC e SI e histórico de 2023 a 2025};
    \node[caixa, below=of f] (e) {Extração, normalização de nomes e vínculo entre PDF e páginas públicas};
    \node[caixa, below=of e] (a) {Classificação curricular, auditoria e tabelas de revisão humana};
    \node[caixa, below=of a] (i) {Instância de 2026 em JSON, com controle de prontidão};
    \node[caixa, right=1.3cm of f] (s) {Perfil sintético: dados observados completados por \textit{proxies} documentadas};
    \node[caixa, below=of s] (g) {Melhoria gulosa integrada a partir da grade observada};
    \node[caixa, below=of g] (m) {SA, ILS e VNS de referência em Python e SA nativo do pyoptframe};
    \node[caixa, below=of m] (v) {Validação estrutural, dupla avaliação, manifestos e relatórios};
    \draw[seta] (f) -- (e);
    \draw[seta] (e) -- (a);
    \draw[seta] (a) -- (i);
    \draw[seta] (i.east) -- ++(0.65,0) |- (s.west);
    \draw[seta] (s) -- (g);
    \draw[seta] (g) -- (m);
    \draw[seta] (m) -- (v);
  \end{tikzpicture}
  \caption{Fluxo de dados e de otimização implementado. Fonte: elaborada pelo autor (2026).}
  \label{fig:pipeline}
\end{figure}

\subsubsection{Instância observada de 2026}

O recorte usado pelo solver reúne as turmas vinculadas às grades de CC e SI nos dois semestres de 2026. A Tabela~\ref{tab:instancia} apresenta sua caracterização. Os 178 registros dos PDFs geraram 180 vínculos com as páginas públicas, porque duas linhas compactadas correspondiam, cada uma, a duas turmas distintas. Todas as 180 turmas têm vagas e inscritos conhecidos, e nenhum encontro ficou sem sala. Das 180 turmas, 179 são do IC e apenas uma pertence a outro departamento. As demais disciplinas externas previstas nas grades, como as de Matemática e Física, não constam do recorte, e por isso ainda não participam da verificação de H8.

\begin{table}[htbp]
  \centering
  \caption{Caracterização do recorte CC/SI da instância observada de 2026. Fonte: elaborada pelo autor a partir dos quadros de horários de 2026/1 e 2026/2 (2026).}
  \label{tab:instancia}
  \small
  \begin{tabular}{lr}
    \toprule
    \textbf{Indicador} & \textbf{Valor} \\
    \midrule
    Turmas físicas & 180 \\
    Encontros semanais & 329 \\
    Docentes observados & 61 \\
    Salas observadas & 22 \\
    Turmas classificadas como obrigatórias & 149 \\
    Turmas com obrigatoriedade ainda indefinida & 4 \\
    Códigos na interseção das grades de CC e SI & 64 \\
    Turmas de 2026 compartilhadas pela interseção curricular & 25 \\
    Turmas com vagas alocadas simultaneamente a CC e SI & 55 \\
    \bottomrule
  \end{tabular}
\end{table}

A auditoria da alocação observada encontrou três sobreposições candidatas de sala, nenhuma de professor e dez conflitos curriculares segundo o contador de H8. Também foi encontrada uma divergência de horário entre o PDF e a página pública, documentada para revisão. Quanto a H12, as 149 turmas obrigatórias comportam no máximo 49 docentes com três obrigatórias cada, sem contar alocações múltiplas. O cumprimento da regra, portanto, depende de qual universo de docentes ela abrange. A tabela de revisão desse universo lista 73 docentes candidatos, ainda sem a indicação oficial de quem está sujeito à regra.

Permanecem pendentes de validação as capacidades físicas das salas, os recursos exigidos por disciplina, os setores e seus dias, as habilitações docentes, as prioridades, o universo de H12, a política de cotutoria e a lista de horários fixos. As capacidades disponíveis são apenas limites inferiores derivados das vagas observadas. \pendente{atualizar esta seção quando as tabelas de revisão forem validadas e aplicadas à instância definitiva.}

\subsubsection{Perfis sintéticos}
\label{subsubsec:perfis}

Para desenvolver e testar o solver sem aguardar todas as validações, a instância observada foi completada por \textit{proxies} determinísticas e documentadas, formando um perfil sintético. Turmas, docentes observados, horários e salas continuam vindo dos dados de 2026; apenas os parâmetros ausentes recebem regras reprodutíveis, resumidas na Tabela~\ref{tab:proxies}. Esses valores não substituem os dados oficiais e servem para validar o software e orientar a etapa definitiva.

\begin{table}[htbp]
  \centering
  \caption{Regras usadas para completar os parâmetros ainda não validados no perfil sintético. Fonte: elaborada pelo autor (2026).}
  \label{tab:proxies}
  \small
  \begin{tabular}{p{4.2cm}p{10.2cm}}
    \toprule
    \textbf{Parâmetro} & \textbf{Regra provisória} \\
    \midrule
    Capacidade da sala & Maior demanda observada na sala em 2025 ou 2026 \\
    Laboratório & Sala cujo identificador começa com a letra L \\
    Recurso por encontro & Inferido do uso histórico de sala comum ou laboratório \\
    Preferência docente & Frequência histórica professor--disciplina, normalizada pela maior frequência do código \\
    Prioridade docente & Valor neutro 1,0 para todos \\
    Setor e habilitação & Setor histórico da disciplina em 2025; habilitados são os professores do setor, incluindo o observado \\
    Universo de H12 & 40 docentes observados, selecionados por carga e por viabilidade conjunta \\
    Cotutoria em H12 & Crédito fracionado entre os professores da turma \\
    Pesos dos critérios fracos & Unitários \\
    \bottomrule
  \end{tabular}
\end{table}

O universo sintético de H12 foi verificado por emparelhamento bipartido entre docentes e turmas obrigatórias. Os domínios cobrem conjuntamente os 118 créditos necessários, o que evita considerar viável um conjunto em que cada docente parece ter opções isoladamente, mas todos disputam as mesmas turmas. O perfil contém 64 docentes no cadastro, 150 turmas obrigatórias do IC e 146 turmas com professor móvel. Cada turma tem até 12 professores candidatos e até 12 padrões de horário.

Foram usadas duas variantes, que diferem apenas no tratamento dos horários:
\begin{itemize}
  \item \textbf{perfil v1}: obrigatórias e turmas de outros departamentos têm horário fixo, e apenas as optativas do IC são flexíveis, o que resulta em 27 turmas com horário variável;
  \item \textbf{perfil v2 (exploratório)}: somente as turmas de outros departamentos e as disciplinas de serviço têm horário fixo, e as obrigatórias do IC podem mudar de faixa dentro dos dias do seu setor, o que resulta em 173 turmas com horário variável e sete fixas.
\end{itemize}

Em ambos os perfis, turmas de projeto final e turmas sem padrão alternativo observado mantêm o horário original. A disciplina de serviço configurada no perfil v2, Estruturas de Dados oferecida a outros cursos, não pertence ao recorte CC/SI; na prática, as sete turmas fixas desse perfil são a única turma externa do recorte e seis turmas do IC sem alternativa de horário.

O perfil v2 materializa duas hipóteses provisórias acordadas pelo autor, ainda pendentes de validação pelo orientador. A primeira restringe H12 aos docentes permanentes do IC, aproximados pelos mesmos 40 docentes do perfil v1. A segunda considera fixos apenas os horários de turmas externas e de serviço. Os domínios de horário de ambos os perfis usam os padrões de dias observados por setor em 2025 e as faixas horárias observadas em 2026. \pendente{registrar a decisão do orientador sobre as hipóteses do perfil v2.}

% =====================================================================
% 2.5 MÉTODO DE SOLUÇÃO
% =====================================================================
\subsection{Método de solução}
\label{subsec:metodo}

\subsubsection{Representação da solução}

A solução é representada diretamente sobre a estrutura JSON da instância. Cada turma guarda o professor atribuído e a lista de encontros, cada um com dia, início, término e sala. Os dados estáticos, como vagas, grupos curriculares e domínios de professores e horários, ficam na mesma estrutura, mas não são alterados pela busca. Indicadores de horário fixo, professor fixo e sala fixa determinam quais decisões cada turma expõe às vizinhanças. O professor observado no quadro original é preservado em um campo separado da atribuição corrente, o que permite comparar a solução com a grade praticada.

\subsubsection{Estruturas de vizinhança}

Foram implementados três tipos de movimento, correspondentes às três decisões do modelo:
\begin{enumerate}
  \item \textbf{professor}: atribui a uma turma do IC outro professor do seu domínio de habilitados; o movimento só se aplica a turmas com um único professor, sem professor fixo e com mais de um candidato;
  \item \textbf{horário}: substitui o padrão de uma turma com horário variável por outro padrão do seu domínio, com o mesmo número de encontros;
  \item \textbf{sala}: troca a sala de um encontro por outra sala compatível em capacidade e em exigência de laboratório.
\end{enumerate}

A cada passo, sorteia-se de maneira uniforme um tipo de movimento entre os disponíveis e, em seguida, a turma, o encontro e o novo valor. Todo movimento devolve uma operação de desfazer que restaura exatamente o estado anterior. Isso evita copiar a solução inteira a cada vizinho avaliado, e a reversibilidade é verificada por testes automatizados para os três tipos. A avaliação, porém, ainda recalcula todos os critérios a cada movimento, sem avaliação incremental.

\subsubsection{Melhoria gulosa integrada}

A busca parte da grade observada, e não de uma solução construída do zero. Uma passagem gulosa percorre as turmas uma vez. Para cada turma, testa até seis professores candidatos, ordenados pela preferência, até seis padrões de horário e até seis salas para cada encontro. Em cada decisão, adota-se o melhor candidato segundo a chave lexicográfica $\kappa$ da Equação~\eqref{eq:chave}, desde que ele melhore a solução corrente. A saída dessa passagem é a solução inicial comum a todas as meta-heurísticas.

\subsubsection{Meta-heurísticas de referência}

As três meta-heurísticas foram implementadas em Python sobre os mesmos movimentos e o mesmo avaliador. Todas recebem a mesma solução inicial e o mesmo orçamento, medido em número de avaliações, e retornam a melhor solução encontrada segundo $\kappa$.

O SA segue o Algoritmo~\ref{alg:sa}. A aceitação usa a regra da Equação~\eqref{eq:metropolis} sobre a energia $E$ definida na Seção~\ref{subsubsec:avaliacao}, e a temperatura decresce geometricamente, $T_{k+1}=\alpha T_k$. Como $E$ multiplica as violações fortes por $10^9$, pioras em restrições fortes são praticamente sempre rejeitadas. Pioras no déficit de H12 só têm probabilidade relevante de aceitação no início da busca, e pioras em critérios fracos são aceitas com mais frequência.

\begin{algorithm}[htbp]
  \DontPrintSemicolon
  \KwIn{solução inicial $s_0$, semente, orçamento $N$, temperatura $T_0$ e fator $\alpha$}
  \KwOut{melhor solução $s^*$ segundo a chave $\kappa$}
  $s\gets s_0$; \quad $s^*\gets s_0$; \quad $T\gets T_0$\;
  \While{o número de avaliações for menor que $N$}{
    sortear um tipo de movimento disponível e um movimento $m$ desse tipo\;
    aplicar $m$ a $s$, obtendo $s'$, e avaliar $s'$\;
    $\Delta\gets E(s')-E(s)$\;
    \eIf{$\Delta\leq 0$ \textbf{ou} $U(0,1)<\exp(-\Delta/T)$}{
      $s\gets s'$\;
      \lIf{$\kappa(s)<\kappa(s^*)$}{$s^*\gets s$}
    }{
      desfazer $m$\;
    }
    $T\gets\alpha T$\;
  }
  \Return{$s^*$}\;
  \caption{\textit{Simulated Annealing} de referência sobre as vizinhanças integradas}
  \label{alg:sa}
\end{algorithm}

O ILS aplica à melhor solução uma perturbação de três movimentos aleatórios e, em seguida, uma busca local aleatória de até 30 tentativas, que só aceita movimentos de melhora. O resultado substitui a melhor solução apenas quando é melhor segundo $\kappa$. O VNS ordena as vizinhanças em professor, horário e sala. Na vizinhança $k$, a agitação aplica $k+1$ movimentos aleatórios daquele tipo, seguida da mesma busca local. Se houver melhora, a busca retorna à primeira vizinhança; caso contrário, passa à seguinte.

\subsubsection{Integração com o pyoptframe}

A integração reaproveita a representação, o avaliador e os movimentos já descritos. Um adaptador registra no motor do pyoptframe a solução, o problema, o construtivo e a função de minimização, além de uma classe única de movimento e uma estrutura de vizinhança. A função de minimização retorna a energia $E$. O construtivo executa a mesma melhoria gulosa, de modo que o motor nativo parte da mesma solução inicial das implementações de referência. Cada movimento aplicado devolve o movimento reverso, como exige o framework.

A busca nativa usa o \textit{BasicSimulatedAnnealing} do OptFrame, com fator de resfriamento $\alpha=0{,}98$, 100 iterações por temperatura e temperatura inicial $T_0=10^8$, calibrada para a escala da energia. O critério de parada do motor é o tempo de execução. A semente controla o sorteio de movimentos feito em Python, mas o gerador de números aleatórios interno do motor não é configurado pelo adaptador. Assim, o fluxo é rastreável, mas a repetição exata de uma execução nativa não é garantida. \pendente{investigar o controle da semente interna do OptFrame e incluir ILS e VNS nativos, se houver tempo.}

\subsubsection{Verificação e rastreabilidade}

Três mecanismos verificam as soluções retornadas. Primeiro, um avaliador de referência, baseado em tabelas, e um avaliador direto sobre o JSON, usado durante a busca, devem produzir os mesmos contadores e critérios. Segundo, uma validação estrutural detecta turmas perdidas ou duplicadas, encontros sem sala ou com sala desconhecida, alterações em dados estáticos e mudanças de professor ou horário fora dos domínios ou em turmas fixas. Terceiro, os executores rejeitam qualquer resultado pior que a solução inicial. Cada execução gera um manifesto com os parâmetros, as sementes e os \textit{hashes} da instância, da configuração e do código. A suíte de testes automatizados do projeto contém 94 testes, todos aprovados na versão usada neste texto.

% =====================================================================
% 2.6 PROTOCOLO EXPERIMENTAL e 2.7 RESULTADOS E DISCUSSÃO
% Números extraídos de:
%   dados/processados/relatorio_experimento_sintetico_2026.md (perfil v1)
%   dados/processados/relatorio_experimento_sintetico_2026_v2_hp.md (perfil v2)
%   dados/processados/piloto_optframe_20260904_final/ (piloto nativo, v1)
%   dados/processados/piloto_optframe_20260907_204501/ (piloto nativo, v2)
%   dados/processados/diagnostico_curricular_2026/ e h8_trajetorias_2026/
% =====================================================================
\subsection{Protocolo experimental}
\label{subsec:protocolo}

Os experimentos realizados até o momento usam os perfis sintéticos v1 e v2 descritos na Seção~\ref{subsubsec:perfis}. Em cada perfil, a grade observada de 2026 é avaliada com as mesmas \textit{proxies} das soluções geradas e serve de referência, chamada aqui de \textit{baseline}. A melhoria gulosa é executada uma vez, e sua saída é a solução inicial de todas as buscas.

SA, ILS e VNS de referência foram executados com cinco sementes (101, 202, 303, 404 e 505) e orçamento de 1.500 avaliações por execução. O SA usou $T_0=10^6$ e $\alpha=0{,}997$. O ILS usou perturbação de três movimentos e busca local de 30 tentativas, e o VNS usou a mesma busca local. No piloto do SA nativo, cada busca teve dez segundos e partiu da mesma solução gulosa, com os parâmetros da seção anterior. As métricas registradas são o total de violações fortes e sua decomposição, o déficit de H12, cada critério fraco, o tempo de execução e o número de avaliações.

As execuções usaram Python 3.14.7 e OptFrame 5.1.0. \pendente{informar processador, memória e sistema operacional da máquina usada.} Os tempos servem apenas como ordem de grandeza, pois as implementações e os critérios de parada diferem entre as buscas de referência e a nativa.

\pendente{definir o protocolo final na instância definitiva: número de execuções por configuração, orçamento comparável entre implementações (avaliações ou tempo), calibração de parâmetros e teste estatístico usado na comparação.}

Também está planejado um estudo da ordem de planejamento das grades, em três cenários. No cenário E1, a grade de CC é otimizada e congelada antes da otimização das decisões ainda livres de SI. O cenário E2 usa a ordem inversa, e o cenário E3 otimiza as duas grades simultaneamente. A comparação entre E1 e E2 mede o custo de planejar uma grade depois da outra, que tende a crescer com o número de turmas compartilhadas. \pendente{executar E1--E3 ou retirá-los do escopo e dos objetivos específicos.}

\subsection{Resultados e discussão}
\label{subsec:resultados}

Os resultados desta seção são preliminares. Eles foram obtidos sobre perfis sintéticos e servem para validar o método, e não para recomendar uma grade à instituição. \pendente{substituir ou complementar esta seção com os resultados da instância definitiva.}

\subsubsection{Perfil v1: horários de obrigatórias fixos}

A Tabela~\ref{tab:resultado_v1} compara o \textit{baseline}, a melhoria gulosa e a melhor execução de cada meta-heurística no perfil v1. A grade observada, avaliada com as \textit{proxies}, apresenta 27 violações fortes: três conflitos de sala, dez conflitos curriculares e 14 docentes do universo sintético abaixo do mínimo anual. A melhoria gulosa reduziu esse total para 12, uma redução de 55,6\%, e também diminuiu dias com aula, janelas e desperdício de capacidade.

\begin{table}[htbp]
  \centering
  \caption{Perfil v1: \textit{baseline}, melhoria gulosa e melhor execução de cada meta-heurística. Fonte: elaborada pelo autor (2026).}
  \label{tab:resultado_v1}
  \small
  \setlength{\tabcolsep}{4pt}
  \begin{tabular}{lrrrrrrrrrr}
    \toprule
    \textbf{Configuração} & \textbf{Hard} & \textbf{Sala} & \textbf{Prof.} & \textbf{Curr.} & \textbf{H12} & \textbf{Dias} & \textbf{Jan.} & \textbf{Desp.} & \textbf{Rod.} & \textbf{Pref.} \\
    \midrule
    \textit{Baseline} & 27 & 3 & 0 & 10 & 14 & 233 & 46 & 3.521 & 38 & 79,28 \\
    Guloso            & 12 & 0 & 0 & 10 &  2 & 215 & 36 & 2.548 & 33 & 76,05 \\
    SA                & 10 & 0 & 0 & 10 &  0 & 230 & 80 & 3.453 & 29 & 68,53 \\
    ILS               & 11 & 0 & 0 & 10 &  1 & 217 & 31 & 2.548 & 32 & 76,05 \\
    VNS               & 11 & 0 & 0 & 10 &  1 & 214 & 29 & 2.544 & 32 & 74,38 \\
    \bottomrule
  \end{tabular}

  \vspace{0.2cm}
  \parbox{0.95\textwidth}{\footnotesize Legenda: Hard = total de violações fortes; Sala, Prof. e Curr. = conflitos de sala, de professor e curriculares; H12 = docentes abaixo do mínimo anual; Dias = combinações professor--semestre--dia com aula; Jan. = janelas; Desp. = desperdício estimado de capacidade; Rod. = repetições professor--disciplina nos dois semestres; Pref. = preferência atendida, ponderada pela prioridade (maior é melhor). Capacidade, recursos e descanso não tiveram violações.}
\end{table}

O SA obteve o menor número de violações fortes. Sua melhor execução, com a semente 303, chegou a 10 violações, uma redução de 62,96\% em relação ao \textit{baseline}, e eliminou os déficits de H12. As dez violações restantes são exatamente os conflitos curriculares do \textit{baseline}. Elas não podem ser corrigidas neste perfil, porque as turmas envolvidas são obrigatórias e têm horário fixo. Para zerar H12, porém, o SA aceitou mais janelas (de 36 na solução gulosa para 80) e reduziu a preferência atendida. ILS e VNS deixaram um docente abaixo do mínimo, mas preservaram melhor os critérios fracos, com menos janelas que o próprio \textit{baseline}.

\subsubsection{Perfil v2: horários de obrigatórias flexíveis}

A Tabela~\ref{tab:resultado_v2} apresenta a mesma comparação no perfil exploratório v2. Com os horários das obrigatórias flexíveis dentro dos dias do setor, a melhoria gulosa reduziu as violações fortes de 27 para 4, uma redução de 85,2\%. A melhor execução do SA, com a semente 505, não apresentou nenhuma violação forte contabilizada: eliminou os conflitos de sala e curriculares e atendeu a carga mínima de todo o universo sintético de H12.

\begin{table}[htbp]
  \centering
  \caption{Perfil v2 (exploratório): \textit{baseline}, melhoria gulosa e melhor execução de cada meta-heurística. Fonte: elaborada pelo autor (2026).}
  \label{tab:resultado_v2}
  \small
  \setlength{\tabcolsep}{4pt}
  \begin{tabular}{lrrrrrrrrrr}
    \toprule
    \textbf{Configuração} & \textbf{Hard} & \textbf{Sala} & \textbf{Prof.} & \textbf{Curr.} & \textbf{H12} & \textbf{Dias} & \textbf{Jan.} & \textbf{Desp.} & \textbf{Rod.} & \textbf{Pref.} \\
    \midrule
    \textit{Baseline} & 27 & 3 & 0 & 10 & 14 & 233 &  46 & 3.521 & 38 & 79,28 \\
    Guloso            &  4 & 0 & 1 &  1 &  2 & 214 &  38 & 2.540 & 32 & 76,55 \\
    SA                &  0 & 0 & 0 &  0 &  0 & 243 & 123 & 5.116 & 25 & 60,60 \\
    ILS               &  2 & 0 & 0 &  1 &  1 & 218 &  38 & 2.582 & 32 & 76,55 \\
    VNS               &  2 & 0 & 0 &  1 &  1 & 216 &  36 & 2.540 & 31 & 75,55 \\
    \bottomrule
  \end{tabular}

  \vspace{0.2cm}
  \parbox{0.95\textwidth}{\footnotesize Legenda: a mesma da Tabela~\ref{tab:resultado_v1}.}
\end{table}

Esse resultado tem custo nos critérios fracos. Em relação ao \textit{baseline}, a solução do SA aumentou as janelas de 46 para 123, o desperdício estimado de capacidade de 3.521 para 5.116 e os dias com aula de 233 para 243. A preferência atendida caiu de 79,28 para 60,60, enquanto o rodízio melhorou, com as repetições caindo de 38 para 25. A solução alterou 89 atribuições de professor, 61 padrões de horário e 196 salas de encontros em relação à grade observada. ILS e VNS terminaram com duas violações, mas mantiveram janelas e desperdício próximos aos da solução gulosa.

A Figura~\ref{fig:hard} resume as violações fortes nos dois perfis. A comparação mostra que a liberdade para mover os horários das obrigatórias é o que permite eliminar os conflitos curriculares. As hipóteses que a sustentam, porém, ainda dependem de validação do orientador.

\begin{figure}[htbp]
  \centering
  \begin{tikzpicture}
    \begin{axis}[
        ybar,
        bar width=14pt,
        width=0.92\textwidth,
        height=6.2cm,
        ymin=0, ymax=31,
        ylabel={Violações fortes},
        symbolic x coords={Baseline,Guloso,SA,ILS,VNS},
        xticklabels={\textit{Baseline},Guloso,SA,ILS,VNS},
        xtick=data,
        enlarge x limits=0.14,
        nodes near coords,
        every node near coord/.append style={font=\footnotesize},
        legend style={at={(0.98,0.97)}, anchor=north east, font=\small},
        tick label style={font=\small},
      ]
      \addplot[fill=gray!35, draw=black] coordinates {(Baseline,27) (Guloso,12) (SA,10) (ILS,11) (VNS,11)};
      \addplot[fill=gray!80, draw=black] coordinates {(Baseline,27) (Guloso,4) (SA,0) (ILS,2) (VNS,2)};
      \legend{Perfil v1, Perfil v2 (exploratório)}
    \end{axis}
  \end{tikzpicture}
  \caption{Violações fortes do \textit{baseline}, da melhoria gulosa e da melhor execução de cada meta-heurística nos perfis sintéticos. Fonte: elaborada pelo autor (2026).}
  \label{fig:hard}
\end{figure}

\subsubsection{Variação entre sementes}

A Tabela~\ref{tab:sementes} resume as cinco execuções de cada método. Nos dois perfis, o SA teve a menor média de violações fortes e foi o único método a atender H12 em alguma execução. No perfil v1, isso ocorreu em três das cinco sementes; no perfil v2, em quatro. ILS e VNS terminaram com pelo menos um docente abaixo do mínimo em todas as execuções. Os resultados por semente estão no Apêndice~C.

\begin{table}[htbp]
  \centering
  \caption{Resumo das cinco execuções por método: média e desvio-padrão amostral das violações fortes, melhor resultado, execuções que atenderam H12, déficit médio de H12 e tempo médio. Fonte: elaborada pelo autor (2026).}
  \label{tab:sementes}
  \small
  \begin{tabular}{llrrrrr}
    \toprule
    & & \multicolumn{2}{c}{\textbf{Violações fortes}} & \textbf{H12} & \textbf{Déficit} & \textbf{Tempo} \\
    \cmidrule(lr){3-4}
    \textbf{Perfil} & \textbf{Método} & \textbf{média $\pm$ dp} & \textbf{melhor} & \textbf{atendida} & \textbf{médio} & \textbf{(s)} \\
    \midrule
    v1 & SA  & 10,4 $\pm$ 0,55 & 10 & 3/5 & 0,4 & 5,8 \\
    v1 & ILS & 11,2 $\pm$ 0,45 & 11 & 0/5 & 1,4 & 8,3 \\
    v1 & VNS & 11,2 $\pm$ 0,45 & 11 & 0/5 & 1,2 & 8,7 \\
    \midrule
    v2 & SA  & 1,0 $\pm$ 0,71 & 0 & 4/5 & 0,2 & 5,2 \\
    v2 & ILS & 2,6 $\pm$ 0,55 & 2 & 0/5 & 1,4 & 5,4 \\
    v2 & VNS & 2,8 $\pm$ 0,84 & 2 & 0/5 & 1,0 & 5,6 \\
    \bottomrule
  \end{tabular}
\end{table}

Com apenas cinco sementes e parâmetros não calibrados, essas diferenças não permitem afirmar superioridade estatística de um método. Ainda assim, o padrão é consistente nos dois perfis. O SA, que aceita pioras nos critérios fracos, alcança a viabilidade com mais frequência. ILS e VNS, que só aceitam melhoras, preservam melhor a qualidade da solução gulosa. \pendente{aplicar teste estatístico com mais execuções na instância definitiva.}

\subsubsection{Piloto com o SA nativo do pyoptframe}

A Tabela~\ref{tab:piloto} compara o SA de referência e o SA nativo, ambos partindo da mesma solução gulosa. Todas as soluções nativas passaram pela validação estrutural e pela conferência entre os dois avaliadores. A reexecução do SA de referência no piloto reproduziu, semente a semente, os mesmos resultados do experimento anterior, o que confirma o determinismo da implementação em Python sob semente fixa.

\begin{table}[htbp]
  \centering
  \caption{Piloto do SA nativo do pyoptframe: violações fortes por semente, na ordem 101, 202, 303, 404 e 505. Fonte: elaborada pelo autor (2026).}
  \label{tab:piloto}
  \small
  \begin{tabular}{llcrr}
    \toprule
    \textbf{Perfil} & \textbf{Implementação} & \textbf{Violações por semente} & \textbf{Melhor} & \textbf{Avaliações} \\
    \midrule
    v1 (guloso: 12) & SA de referência & 11, 10, 10, 10, 11 & 10 & 1.500 \\
                    & SA nativo        & 12, 12, 10, 12, 11 & 10 & 422 a 471 \\
    \midrule
    v2 (guloso: 4)  & SA de referência & 2, 1, 1, 1, 0 & 0 & 1.500 \\
                    & SA nativo        & 2, 3, 4, 3, 4 & 2 & 341 a 369 \\
    \bottomrule
  \end{tabular}
\end{table}

No perfil v1, o SA nativo alcançou o mesmo melhor resultado da referência em uma das cinco sementes, mas três execuções terminaram sem melhorar a solução gulosa. No perfil v2, que tem um espaço de busca maior, ele ficou atrás da referência em todas as sementes, e duas execuções também não melhoraram a solução gulosa. Os dois arranjos, porém, não são diretamente comparáveis. O SA nativo teve dez segundos por execução e realizou de 341 a 471 avaliações, enquanto a referência usou 1.500 avaliações em cerca de cinco a seis segundos. Os esquemas de resfriamento e as temperaturas também diferem. Uma hipótese para a menor vazão do motor nativo é o custo da ponte entre Python e o núcleo do OptFrame, que ainda precisa ser medido. O piloto, portanto, valida a integração, mas não estabelece qual implementação é melhor.

\subsubsection{Conflitos curriculares e a interpretação de H8}
\label{subsubsec:h8}

Os dez conflitos curriculares da grade observada foram diagnosticados individualmente. Eles ocorrem em dois grupos da grade de SI: quatro sobreposições no primeiro período em cada semestre e duas no sexto período em 2026/1. Os horários envolvidos coincidem nas extrações dos PDFs e na coleta das páginas públicas, portanto as sobreposições são reais. Em todos os casos, porém, há uma escolha de uma seção por disciplina sem choque. Em 2026/1, por exemplo, a turma TCC00332-A1 coincide com TCC00354-A1, mas a combinação de TCC00332-A1 com TCC00354-B1 evita o choque.

Isso expõe uma questão de modelagem. O contador atual de H8 exige que nenhuma oferta de disciplinas diferentes do grupo se sobreponha, mesmo quando essas ofertas são seções alternativas. Uma interpretação alternativa exigiria apenas que o aluno tivesse ao menos uma trajetória viável, isto é, uma escolha de seções sem choque. Essa segunda leitura é mais próxima da matrícula real, mas exige modelar a elegibilidade das seções e, para garantir atendimento coletivo, a demanda e as vagas.

Para apoiar essa decisão, foi implementada uma verificação paralela de H8 por trajetória, que não altera o avaliador usado na busca. Para cada grupo curricular, ela procura uma seção utilizável de cada disciplina obrigatória do período, sem choques entre as escolhidas. Como filtro provisório de elegibilidade, considera-se utilizável a seção com vagas ofertadas para o curso do grupo. A Tabela~\ref{tab:trajetorias} mostra o resultado para os 32 grupos da grade observada, que correspondem a oito períodos de cada curso nos dois semestres.

\begin{table}[htbp]
  \centering
  \caption{Verificação de H8 por trajetória nos 32 grupos curriculares de 2026. Fonte: elaborada pelo autor (2026).}
  \label{tab:trajetorias}
  \small
  \begin{tabular}{p{6cm}rrr}
    \toprule
    & \multicolumn{3}{c}{\textbf{Grupos curriculares}} \\
    \cmidrule(lr){2-4}
    \textbf{Dados considerados} & \textbf{com trajetória} & \textbf{sem trajetória} & \textbf{incompletos} \\
    \midrule
    Somente o recorte CC/SI & 7 & 0 & 25 \\
    Recorte com complemento da coleta web & 22 & 0 & 10 \\
    \bottomrule
  \end{tabular}
\end{table}

No recorte, a maioria dos grupos fica inconclusiva porque as disciplinas de outros departamentos não estão na instância. Com o complemento da coleta web, que só acrescenta disciplinas ausentes e nunca altera horários da solução, 22 grupos têm trajetória e nenhum grupo fica comprovadamente sem trajetória. Os dez grupos restantes têm, cada um, uma disciplina sem turma utilizável nos dados. Esse diagnóstico não comprova vagas suficientes nem atendimento de todos os alunos, mas indica que, nos grupos afetados, os conflitos contados por H8 na grade observada não impedem que um aluno curse o período. \pendente{registrar a decisão do orientador sobre a semântica de H8 e, se for o caso, incorporar a verificação por trajetória ao avaliador.}

\subsubsection{Discussão}

Três observações emergem dos resultados preliminares. A primeira é que a melhoria gulosa responde pela maior parte do ganho: nos dois perfis, ela removeu mais da metade das violações fortes com uma única passagem. A segunda é que os métodos se complementam. O SA é mais eficaz em alcançar viabilidade, enquanto ILS e VNS preservam melhor os critérios fracos, o que sugere um fluxo combinado em que o SA busca a viabilidade e ILS ou VNS refinam a solução.

A terceira observação é que existe um conflito claro entre viabilidade e conforto docente. Como a chave lexicográfica prioriza qualquer redução de violações fortes, soluções viáveis podem ter mais janelas, mais dias de aula e menos preferências atendidas. Esse comportamento é coerente com a formulação, mas reforça a necessidade de definir com o orientador quais regras são de fato inegociáveis e quais pesos devem equilibrar os critérios fracos.

Os resultados têm limitações importantes. Capacidades, recursos, habilitações, prioridades e o universo de H12 são \textit{proxies}, e as capacidades são apenas limites inferiores. O perfil v2 depende de hipóteses ainda não validadas. A instância contém apenas uma turma de outro departamento, de modo que H8 ainda não considera a maior parte das disciplinas externas das grades. Por fim, o número de execuções é pequeno e os orçamentos das implementações de referência e nativa não são equivalentes. \pendente{revisar esta discussão após os experimentos definitivos, incluindo gráficos de convergência e a comparação com a grade observada na instância oficial.}


% =====================================================================
% 3. CONCLUSÃO
% =====================================================================
\section{Conclusão}
\label{sec:conclusao}

Este trabalho abordou a alocação anual de professores, horários e salas no IC/UFF por meio de uma abordagem meta-heurística implementada com o pyoptframe. O problema foi formalizado como uma variante do CB-CTT que incorpora a atribuição de professores, o horizonte anual, a sala por encontro, as grades de CC e SI com turmas compartilhadas e regras locais, como setores, horários fixos, descanso entre jornadas e a carga mínima anual de obrigatórias.

Uma pipeline reprodutível transformou os quadros de horários de 2026, as páginas públicas de oferta e as grades curriculares em uma instância auditável, com 180 turmas, 329 encontros semanais e 22 salas. A pipeline separa os dados observados, as decisões institucionais pendentes e as hipóteses sintéticas, o que impede que uma \textit{proxy} seja tratada como dado oficial. Sobre essa base, foram implementados um avaliador decomposto por critério, movimentos de professor, horário e sala, uma melhoria gulosa e as meta-heurísticas SA, ILS e VNS, além da integração com o SA nativo do pyoptframe. As soluções são conferidas por dois avaliadores e por uma validação estrutural.

Os resultados preliminares, obtidos em perfis sintéticos, mostram que a busca reduz substancialmente as violações de restrições fortes da grade observada. Com os horários das obrigatórias fixos, a melhor solução passou de 27 para 10 violações, todas curriculares e impossíveis de corrigir nesse perfil. Com horários flexíveis dentro dos setores, cenário que depende de hipóteses ainda não validadas, o SA obteve uma solução sem violações fortes contabilizadas. Em ambos os casos, a viabilidade custou mais janelas, mais dias de aula e menos preferências atendidas, o que evidencia o conflito entre os papéis institucionais descrito na introdução. O SA foi o método mais eficaz para alcançar a viabilidade, enquanto ILS e VNS preservaram melhor os critérios fracos.

O trabalho também produziu um resultado de modelagem. Os dez conflitos curriculares da grade observada envolvem seções alternativas, e existe, em cada grupo afetado, uma escolha de seções sem choque. A forma de tratar essas alternativas em H8 é, portanto, uma decisão de modelagem que afeta diretamente a viabilidade das instâncias. \pendente{atualizar este parágrafo com a decisão do orientador sobre H8.}

As principais limitações são o uso de \textit{proxies} para parâmetros institucionais, o pequeno número de execuções, a ausência de calibração de parâmetros e a comparação ainda não equivalente entre as implementações de referência e nativa. Além disso, a instância ainda não inclui a maior parte das disciplinas externas das grades. \pendente{revisar as conclusões após os experimentos na instância definitiva.}

\subsection{Trabalhos futuros}
\label{subsec:trabalhos_futuros}

Como continuidade, destacam-se:
\begin{itemize}
  \item aplicar à instância as decisões institucionais registradas nas tabelas de revisão e repetir os experimentos sem \textit{proxies};
  \item definir a semântica de H8 para turmas alternativas e incluir as disciplinas externas das grades;
  \item estabelecer um protocolo com orçamento comparável entre as implementações, controle da semente interna do OptFrame, mais execuções e teste estatístico;
  \item executar os cenários E1, E2 e E3 para medir o efeito da ordem de planejamento das grades de CC e SI;
  \item implementar avaliação incremental, novas vizinhanças, como a troca de professores e de salas entre turmas, e o fluxo combinado em que o SA busca viabilidade e ILS ou VNS refinam a solução;
  \item formular um modelo de programação inteira para instâncias reduzidas, a fim de obter limitantes e validar as soluções heurísticas;
  \item desenvolver uma forma de visualização das grades geradas, para apoiar a discussão com a chefia e as coordenações.
\end{itemize}


% =====================================================================
% ELEMENTOS PÓS-TEXTUAIS
% =====================================================================
\begin{singlespace}
\bibliography{referencias}
\end{singlespace}

% =====================================================================
% Elementos pós-textuais, na ordem do modelo do curso:
% Apêndice A (lista de ilustrações), Apêndice B (lista de tabelas),
% apêndices do autor, Anexos (opcional) e Agradecimentos (opcional).
% =====================================================================

% ---------------------------------------------------------------------
\titulopostextual{APÊNDICE A - Lista de ilustrações}
\begin{singlespace}
\makeatletter\@starttoc{lof}\makeatother
\end{singlespace}

% ---------------------------------------------------------------------
\titulopostextual{APÊNDICE B - Lista de tabelas}
\begin{singlespace}
\makeatletter\@starttoc{lot}\makeatother
\end{singlespace}

% ---------------------------------------------------------------------
\titulopostextual{APÊNDICE C - Resultados por semente}

A Tabela~\ref{tab:apendice_sementes} apresenta o resultado de cada execução das meta-heurísticas de referência nos perfis sintéticos v1 e v2, com orçamento de 1.500 avaliações. O déficit de H12 é o número de obrigatórias que ainda faltam para que todos os docentes do universo sintético atinjam o mínimo anual.

\begin{table}[htbp]
  \centering
  \caption{Violações fortes, déficit de H12 e tempo de cada execução. Fonte: elaborada pelo autor (2026).}
  \label{tab:apendice_sementes}
  \small
  \begin{tabular}{llrrrrrr}
    \toprule
    & & \multicolumn{3}{c}{\textbf{Perfil v1}} & \multicolumn{3}{c}{\textbf{Perfil v2}} \\
    \cmidrule(lr){3-5}\cmidrule(lr){6-8}
    \textbf{Método} & \textbf{Semente} & \textbf{Hard} & \textbf{Déficit} & \textbf{Tempo (s)} & \textbf{Hard} & \textbf{Déficit} & \textbf{Tempo (s)} \\
    \midrule
    SA & 101 & 11 & 1 & 6,5 & 2 & 0 & 5,3 \\
     & 202 & 10 & 0 & 5,1 & 1 & 0 & 5,1 \\
     & 303 & 10 & 0 & 5,4 & 1 & 0 & 4,9 \\
     & 404 & 10 & 0 & 5,6 & 1 & 1 & 5,0 \\
     & 505 & 11 & 1 & 6,3 & 0 & 0 & 5,7 \\
    \midrule
    ILS & 101 & 11 & 1 & 6,9 & 2 & 1 & 5,3 \\
     & 202 & 11 & 1 & 7,8 & 2 & 1 & 5,5 \\
     & 303 & 11 & 1 & 8,4 & 3 & 2 & 5,4 \\
     & 404 & 11 & 2 & 9,0 & 3 & 1 & 5,4 \\
     & 505 & 12 & 2 & 9,5 & 3 & 2 & 5,6 \\
    \midrule
    VNS & 101 & 11 & 1 & 9,4 & 3 & 1 & 5,7 \\
     & 202 & 11 & 1 & 8,8 & 2 & 1 & 5,7 \\
     & 303 & 11 & 1 & 8,9 & 2 & 1 & 5,7 \\
     & 404 & 11 & 1 & 8,9 & 3 & 1 & 5,7 \\
     & 505 & 12 & 2 & 7,5 & 4 & 1 & 5,4 \\
    \bottomrule
  \end{tabular}
\end{table}

% ---------------------------------------------------------------------
\titulopostextual{APÊNDICE D - Reprodução dos experimentos}

Os dados, o código e os relatórios deste trabalho estão organizados em um repositório próprio. \pendente{informar o endereço público do repositório, se ele for disponibilizado.} A partir da raiz do repositório, os resultados dos perfis sintéticos podem ser regenerados com os comandos abaixo. Cada execução registra, em um manifesto, os parâmetros, as sementes e os \textit{hashes} da instância, da configuração e do código.

\begin{singlespace}
\small
\begin{verbatim}
python scripts/run_pipeline_2026.py --offline
python scripts/build_synthetic_instance_2026.py
python scripts/run_synthetic_experiments_2026.py
python scripts/benchmark_optframe_sa_2026.py 10
python -m unittest discover -s tests
\end{verbatim}
\end{singlespace}

O perfil v2 usa a configuração \texttt{dados/config\_sintetica\_2026\_v2\_hp.json}, informada aos executores pelas opções \texttt{-{}-config} e \texttt{-{}-instance}. \pendente{conferir os comandos exatos do perfil v2 antes da versão final.}

% ---------------------------------------------------------------------
% ANEXO A (opcional): documentos de autoria externa. Descomente se usar.
% \titulopostextual{ANEXO A - Título do anexo}

% ---------------------------------------------------------------------
\titulopostextual{Agradecimentos}

\pendente{redigir os agradecimentos pessoais ao orientador, à família e a quem colaborou com o trabalho.}

Ferramentas de inteligência artificial generativa foram usadas como apoio à programação, à organização de arquivos, à clarificação textual e à discussão do trabalho, incluindo a proposição de alternativas e a análise de modelagem e de métodos experimentais. Hipóteses provisórias foram identificadas como tais e não foram apresentadas como diretrizes confirmadas. As decisões finais permanecem sob responsabilidade do autor, que revisou todo o conteúdo antes de incorporá-lo ao trabalho. \pendente{revisar esta declaração para que ela descreva com precisão o uso de IA na redação deste artigo, cuja primeira versão foi redigida com apoio de um assistente de IA a partir dos registros do projeto.}


\end{document}
